\section{Ullrich 2024: inequalities between s-numbers}
\label{sec:ullrich-2024-s-numbers}

\paragraph{Bibliographic record.}
Mario Ullrich, ``Inequalities between s-numbers,''
\emph{Advances in Operator Theory} 9 (2024), article 82,
DOI \texttt{10.1007/s43036-024-00386-x}.
The article is open access.

\paragraph{Role in this repository.}
This paper is used as a modern theorem-level reference for the axiomatic
s-number interface underlying the approximation-number and finite Ky Fan
layers.  It is not a substitute for the Davis--Kahan source proof.  Its role is
to identify which properties are elementary consequences of the best
finite-rank approximation definition and which are substantive operator-ideal
theorems.

\paragraph{Approximation numbers.}
For a bounded operator $T:X\to Y$, the approximation numbers are the distances
from $T$ to operators of bounded finite rank.  The repository uses a zero-based
index convention:
\[
  a_n(T)=\inf\{\lVert T-R\rVert:\operatorname{rank}R\le n\}.
\]
Consequently $a_0(T)=\lVert T\rVert$.  This differs by one from references that
start with $a_1$ and use rank strictly less than the index.

\paragraph{Axioms used by the formalization.}
The relevant s-number laws are monotonicity, the rank-perturbation/addition
inequality, two-sided ideal control, normalization at the first index, and
vanishing beyond finite rank.  The approximation-number definition gives the
order, normalization, perturbation, homogeneity, and composition laws through
finite-rank competitors.  Adjoint invariance is Hilbert-space-specific and
requires a separate rank-preservation argument.

\paragraph{Ky Fan passage.}
The finite Ky Fan gauge is
\[
  \lVert T\rVert_{(k)}=\sum_{j=0}^{k-1} a_j(T).
\]
Its triangle inequality is not obtained by summing the fixed-index Lipschitz
bound.  It uses the stronger rank-additive approximation inequality and the
associated weak-majorization argument.  The Lean development therefore keeps
this as a named theorem rather than hiding it inside a concrete norm-family
constructor.

\paragraph{Formalization map.}
The best-finite-rank definition and elementary laws live in
\texttt{ForMathlib/Analysis/Normed/Operator/ApproximationNumber.lean}.
The Hilbert-space adjoint law, Ky Fan triangle theorem, cutoff convergence, and
Fan-dominant ideal-family construction live in
\texttt{DavisKahan/Experimental/InfiniteDimensional/Ideals/ApproximationNumbers.lean}.
